\documentclass[UTF8,a4paper,11pt,openany]{ctexbook}
\usepackage{../common/statlearnbook}
\SLSetBookMetadata{第 5 册}{统计学习方法讲义 v2.6.0}{采样方法、主题模型与图排序}
\begin{document}
\setcounter{page}{761}
\setcounter{chapter}{32}
\setcounter{figure}{3}
\pagestyle{headings}
\chaptermark{马尔可夫链蒙特卡罗法}
\sectionmark{平稳分布与细致平衡}
\begin{definition}[细致平衡]
若对任意可测集合$A,B\in\mathcal B$都有
\begin{equation}
\int_A\pi(\mathrm dx)K(x,B)=\int_B\pi(\mathrm dy)K(y,A),
\end{equation}
则称$K$关于$\pi$满足细致平衡，或称平稳链关于$\pi$可逆。若$K$有密度$k(x,y)$，条件化为
\[
\pi(x)k(x,y)=\pi(y)k(y,x)\quad\text{对}\ \mu\otimes\mu\ \text{几乎处处成立}.
\]
\end{definition}
从局部流到平稳分布的推导分三步。第一步，密度形式把$x\to y$的平稳概率流与$y\to x$的反向流配对。第二步，对$x$积分；第三步使用$\int k(y,x)\mu(\mathrm dx)=1$，于是总流入等于目标质量。图\ref{fig:V5-C03-mh-balance-flux}把成对流与接受机制分开绘制。
\input{../../绘图源码/第05册_采样方法主题模型与图排序/V5-C03/fig_v5_c03_mh_balance_flux.tex}
\begin{proposition}[细致平衡推出平稳]
若概率测度$\pi$与马尔可夫核$K$满足细致平衡，则$\pi K=\pi$。
\end{proposition}
\end{document}
